% ===== Section 3: why people like being scared (~18 min) =====
\partaccent{peach}
\topicopener{pumpkin}{Սարսափ}{Մղձավանջներ մի՛ տես... եթե, իհարկե, դա քո սիրածը չէ}{Ինչու են մարդիկ սիրում վախենալ ու իրենք են դա փնտրում}{smug}

\begin{frame}{Վախենալ՝ կամովին}
  \begin{columns}[c]
    \begin{column}{0.46\textwidth}
      \begin{poll}
        Բարձրացրու ձեռքդ, եթե վերջին մեկ տարում կամովին.\par\smallskip
        \emoji{film}\ սարսափ ֆիլմ ես նայել\par
        \emoji{rollercoaster}\ ամերիկյան բլուրներ ես նստել\par
        \emoji{ghost}\ սարսափ պատմություն ես լսել
      \end{poll}
    \end{column}
    \begin{column}{0.48\textwidth}
      \onslide<2->{%
      \begin{card}[colback=peach!40!cardwhite]
        Մինչ այժմ ուղեղն ամեն կերպ \textbf{պաշտպանում} էր մեզ վտանգից. նյարդայնություն, խուսափում, ֆոբիաներ։
      \end{card}}
      \onslide<3->{%
      \begin{card}[colback=brainpink!30!cardwhite]
        Ուրեմն ինչու՞ են մարդիկ կամովին ցատկում կամրջից ռետինե պարանով (bungee) կամ ինքնաթիռից՝ պարաշյուտով \emoji{exploding}
      \end{card}}
    \end{column}
  \end{columns}
\end{frame}

\begin{frame}{Վախը մեծ բիզնես է}
  \begin{columns}[c]
    \begin{column}{0.56\textwidth}
      \begin{itemize}
        \item \emoji{books}\ Սթիվեն Քինգն ու Դին Քունցը միասին մոտ \textbf{մեկ միլիարդ} գիրք են վաճառել
        \item<2-> \emoji{film}\ «Սղոց» (Saw). արյունոտ ֆիլմերի շարք։ Գրքի պահին՝ 7 ֆիլմ, այսօր՝ 10
        \item<3-> \emoji{ghost}\ Սարսափ պատմություններ խարույկի մոտ, «սարսափների տներ»
        \item<4-> \emoji{pumpkin}\ Հելոուինին մեռելի պես ենք հագնվում, որ հարևաններից կոնֆետ «կորզենք»
      \end{itemize}
    \end{column}
    \begin{column}{0.38\textwidth}
      \onslide<2->{%
      \begin{bookquote}
        Բոլորը ցուցադրվել են կինոթատրոններում, այլ ոչ թե կապարե արկղերում փակվել ու ուղարկվել դեպի Արև։
      \end{bookquote}}
      \onslide<5->{%
      \begin{card}[colback=peach!40!cardwhite]
        Ինչու՞ ենք սիրում զվարճանքներ, որոնց իմաստը \textbf{վախեցնելն} է
      \end{card}}
    \end{column}
  \end{columns}
  \only<2->{\source{Saw ֆիլմերի քանակը՝ 2023-ի «Saw X»-ով}}
\end{frame}

% The mesolimbic reward pathway
\begin{frame}{Պարգևի ուղին}
  \begin{columns}[c]
    \begin{column}{0.5\textwidth}
      \centering
      \begin{tikzpicture}[scale=0.9]
        \brainsilhouette
        \node[tag, font=\blackfont\scriptsize] at (-2.6,2.2) {ճակատ};
        \filldraw[fill=amOrange, draw=plum, line width=0.8pt] (0.55,-0.6) circle (0.2);
        \filldraw[fill=amRed, draw=plum, line width=0.8pt] (-1.15,-0.15) circle (0.2);
        \draw[-{Stealth[length=2.4mm]}, amRed, line width=1.6pt] (0.35,-0.55) to[bend left=12] (-0.93,-0.2);
        \node[tag, font=\scriptsize\bfseries] at (-0.25,0.2) {դոֆամին};
        \node[tag, anchor=north, font=\scriptsize, text width=3cm] (vta) at (1.4,-2.4) {փորային ծածկի\\ շրջան (VTA)};
        \draw[plum, line width=0.6pt] (vta.north) -- (0.6,-0.8);
        \node[tag, anchor=north, font=\scriptsize, text width=2.6cm] (nac) at (-2.0,-1.35) {կից կորիզ\\ (nucleus accumbens)};
        \draw[plum, line width=0.6pt] (nac.north) -- (-1.2,-0.35);
      \end{tikzpicture}
    \end{column}
    \begin{column}{0.46\textwidth}
      \textbf{Մեզոլիմբիկ պարգևի ուղի} (mesolimbic reward pathway)։ Վախի սարսուռն ու քաղցրի հաճույքը, հավանաբար, այստեղից են։\par\smallskip
      \begin{itemize}
        \item<2-> \textcolor{amOrange}{$\bullet$} \textbf{VTA}՝ որոշում է՝ լավ է, թե վատ
        \item<3-> \textcolor{amRed}{$\bullet$} \textbf{Կից կորիզ}՝ ստիպում է զգալ դա
        \item<4-> \emoji{doughnut}\ Համեղ նախուտեստ $\rightarrow$ հաճույք
        \item<5-> \emoji{vomit}\ Փչացած կաթ $\rightarrow$ զզվանք. «Էլ երբե՛ք»
      \end{itemize}
    \end{column}
  \end{columns}
  \source{Koob, Trends Pharmacol. Sci., 1992. Նկարը սխեմատիկ է}
\end{frame}

% What counts as a reward; the rat and the PlayStation
\begin{frame}{Ի՞նչ է պարգևը}
  \begin{columns}[c]
    \begin{column}{0.5\textwidth}
      Հաճելի զգացողություն, երբ անում ենք մի բան, որին ուղեղը \textbf{հավանություն} է տալիս.\par\smallskip
      \begin{itemize}
        \item \emoji{bread}\ ուտելիք, հատկապես ածխաջրեր (խեղճ դիետա պահողներ)
        \item<2-> \emoji{kiss}\ սեքս. թեև առանց դրա կարող ենք ապրել։ Այո՛, կարող ենք
        \item<3-> \emoji{hand}\ քոր եկած տեղը քորելը
      \end{itemize}
    \end{column}
    \begin{column}{0.46\textwidth}
      \centering
      \onslide<4->{%
      \begin{tikzpicture}
        \node at (0,0) {\bigemoji[1.1cm]{rat}};
        \draw[plum, line width=2pt] (0.8,-0.45) -- (1.1,0.45); \fill[plum] (1.1,0.45) circle (0.12);
        \node[pbox=mint!40!cardwhite, text width=4.2cm, align=left] at (1.4,-1.4) {\emoji{apple}\ Միրգ $\rightarrow$ ավելի հաճախ է սեղմում};
        \node[pbox=brainpink!30!cardwhite, text width=4.2cm, align=left, visible on=<5->] at (1.4,-2.45) {\emoji{videogame}\ Նոր PlayStation խաղ $\rightarrow$ անտարբեր է};
      \end{tikzpicture}\par\smallskip}
      \onslide<5->{\small Դեռահասը կվիճեր, բայց առնետի համար դա պարգև չէ։}
    \end{column}
  \end{columns}
  \vspace{1mm}
  \onslide<6->{\centering Պարգևը տարբեր է տարբեր մարդկանց համար. ոմանց համար՝ նաև \textbf{վախը}։\par}
  \source{Reyes-Castro et al., 2011}
\end{frame}

% Route 1: curiosity
\begin{frame}{Ճանապարհ 1. հետաքրքրասիրություն}
  \begin{columns}[c]
    \begin{column}{0.58\textwidth}
      \begin{itemize}
        \item Նույնիսկ առնետներն են հետազոտում նոր առարկաները
        \item<2-> Մարդիկ՝ առավել ևս. մեզ ձգում է \textbf{նորությունը}
        \item<3-> Երեխա ունեցողները լավ գիտեն «տեսնենք՝ ինչ կլինի» մղումը։ Հաճախ ավերիչ
      \end{itemize}
      \vspace{2mm}
      \onslide<4->{%
      \begin{card}[colback=peach!40!cardwhite]
        Բայց նոր, անվտանգ բաներն էլ անթիվ են։ Ինչու՞ ընտրել հենց \textbf{վախն ու վտանգը}։
      \end{card}}
    \end{column}
    \begin{column}{0.36\textwidth}
      \centering
      \bigemoji[1.2cm]{baby}\ \bigemoji[0.9cm]{wastebasket}\par\medskip
      \mascot[2.2cm]{thinking}
    \end{column}
  \end{columns}
  \source{Sluckin, Hargreaves \& Colman, 1983}
\end{frame}

% Route 2: relief when the danger stops (schematic curve)
\begin{frame}{Ճանապարհ 2. երբ վատը դադարում է}
  \begin{columns}[c]
    \begin{column}{0.5\textwidth}
      \centering
      \begin{tikzpicture}
        \draw[-{Stealth[length=2.4mm]}, plum, line width=1pt] (0,0) -- (5.6,0) node[below left, tag, font=\small] {ժամանակ};
        \draw[-{Stealth[length=2.4mm]}, plum, line width=1pt] (0,0) -- (0,3.3);
        \draw[amBlue, line width=2pt] (0.1,0.3) .. controls (1.0,0.4) and (1.4,2.8) .. (2.2,2.8) .. controls (2.7,2.8) and (2.8,0.4) .. (3.4,0.3) -- (5.3,0.3);
        \node[tag=amBlue, font=\small\bfseries, anchor=south] at (1.95,2.85) {վախ};
        \draw[amOrange, line width=2pt, dashed, visible on=<2->] (2.9,0.2) .. controls (3.2,0.3) and (3.3,2.3) .. (3.7,2.3) .. controls (4.1,2.3) and (4.3,0.6) .. (5.3,0.5);
        \node[tag=peachdark, font=\small\bfseries, anchor=south, visible on=<2->] at (3.9,2.35) {պարգև};
        \node[tag, font=\scriptsize, text width=2cm, visible on=<2->] at (2.5,-0.4) {վտանգն անցավ};
      \end{tikzpicture}\par
      {\scriptsize Սխեմատիկ, ոչ թե չափումներ}
    \end{column}
    \begin{column}{0.46\textwidth}
      \begin{itemize}
        \item \emoji{hauntedhouse}\ Սարսափների տանն էիր, հիմա դրսում ես
        \item<2-> \emoji{airplane}\ Ընկնում էիր, հիմա գետնին ես ու ողջ
        \item<3-> \emoji{book}\ Պատմությունն ավարտվեց, իսկ մարդասպանը չհայտնվեց
      \end{itemize}
      \vspace{1mm}
      \onslide<4->{%
      \begin{card}[colback=peach!40!cardwhite]
        Ուտելը կյանքդ մի քիչ է բարելավում։ Իսկ այստեղ դու \textbf{մահից խուսափեցիր}. պարգևը շատ ավելի ուժեղ է։
      \end{card}}
    \end{column}
  \end{columns}
\end{frame}

% Vivid memories and feeling alive
\begin{frame}{Ավելի կենդանի, քան երբևէ}
  \begin{columns}[c]
    \begin{column}{0.56\textwidth}
      \begin{itemize}
        \item Ադրենալինից հետո ամեն ինչ սուր է զգացվում, իսկ թեթևացումը շատ ուժեղ է
        \item<2-> Պարգևի ուղին կապված է \textbf{հիպոկամպի} ու \textbf{նշաձև մարմնի} հետ. դեպքը հիշվում է հատկապես վառ
        \item<3-> Վախի պահին կարող ես քեզ ավելի «կենդանի» զգալ, քան երբևէ, իսկ մնացածը թվում է խամրած
      \end{itemize}
    \end{column}
    \begin{column}{0.38\textwidth}
      \centering
      \onslide<4->{%
      \bigemoji[1.1cm]{coffee}\par\smallskip
      \begin{bookquote}
        Ինչպես կրկնակի էսպրեսոյին սովոր մարդը չի բավարարվի կաթով լի լատեով։
      \end{bookquote}}
    \end{column}
  \end{columns}
  \source{Wittmann et al., Neuropsychologia, 2008}
\end{frame}

% It has to be real: video games, fans and harnesses
\begin{frame}{Պետք է իրական լինի}
  \begin{columns}[T]
    \begin{column}{0.47\textwidth}
      \begin{card}
        \emoji[1.4em]{videogame}\ {\blackfont Տեսախաղ}\par\smallskip
        Որքան էլ իրատեսական լինի, մրցարշավի խաղը չի տա իրական արագության սարսուռը։ Ուղեղը տարբերում է իրականն անիրականից։
      \end{card}
    \end{column}
    \begin{column}{0.47\textwidth}
      \onslide<2->{%
      \begin{card}
        \emoji[1.4em]{rollercoaster}\ {\blackfont Էքստրեմալ սպորտ}\par\smallskip
        Էկրաններով ու հսկա օդափոխիչներով bungee-ն չես նմանակի։ Դրա համար էլ կան ամերիկյան բլուրները։
      \end{card}}
    \end{column}
  \end{columns}
  \vspace{3mm}
  \onslide<3->{%
  \begin{card}[colback=peach!40!cardwhite]
    \centering Իսկական սարսուռի համար պետք է տհաճ հետևանքի \textbf{իրական} հնարավորություն։ Ինչ էլ ասեն «տեսախաղերը բռնություն են ծնում» պնդողները։
  \end{card}}
\end{frame}

% Control: why books and films scare more than games
\begin{frame}{Իսկ ինչու՞ են գրքերն ու ֆիլմերը վախեցնում}
  \begin{columns}[T]
    \begin{column}{0.47\textwidth}
      \begin{card}[colback=mint!35!cardwhite]
        \emoji[1.4em]{videogame}\ {\blackfont Խաղում՝ ղեկը քոնն է}\par
        Կարող ես կանգնեցնել, աշխարհը պատասխանում է քեզ
      \end{card}
    \end{column}
    \begin{column}{0.47\textwidth}
      \onslide<2->{%
      \begin{card}[colback=brainpink!30!cardwhite]
        \emoji[1.4em]{film}\emoji[1.4em]{book}\ {\blackfont Ֆիլմում ու գրքում՝ ոչ}\par
        Գիրքը կփակես, բայց պատմությունը չի փոխվի
      \end{card}}
    \end{column}
  \end{columns}
  \vspace{2mm}
  \begin{itemize}
    \item<3-> Որքան շատ է ուղեղը վերահսկում, այնքան \textbf{պակաս} է վախենում
    \item<4-> Վառ հիշողությունները մնում են. ֆիլմը վախեցնում է դեռ օրեր անց
    \item<5-> «Երևակայությանը թողնվածը» կարող է ավելի սարսափելի լինել, քան ամենաարյունոտ էֆեկտները
  \end{itemize}
\end{frame}

% 1970s suggestion vs CGI; survival horror; chat
\begin{frame}{Ակնարկն ավելի վախեցնող է, քան արյունը}
  \begin{columns}[c]
    \begin{column}{0.58\textwidth}
      \begin{itemize}
        \item 1970-ականները սարսափի «ոսկե դարն» են. CGI չկար, վախը գալիս էր ակնարկից ու մթնոլորտից
        \item<2-> Ուղեղն ինքն էր վտանգ «կանխատեսում», ու մարդիկ ցատկում էին ստվերներից
        \item<3-> Երբ ամեն ինչ ուղիղ ցույց են տալիս, ուղեղը հիշում է, որ սա հորինվածք է
        \item<4-> Խաղերում՝ «գոյատևման սարսափ» (survival horror). պետք է \textbf{խուսափել} վտանգից, ոչ թե լազերային թնդանոթով այն փոշիացնել
      \end{itemize}
    \end{column}
    \begin{column}{0.36\textwidth}
      \onslide<5->{%
      \begin{chatbox}
        Ամենասարսափելի ֆիլմը, որ տեսել ես։ Ի՞նչն էր վախեցնում. այն, ինչ \textbf{տեսար}, թե՞ այն, ինչ \textbf{չտեսար}
      \end{chatbox}}
      \centering\bigemoji[1.0cm]{popcorn}
    \end{column}
  \end{columns}
  \source{Tinwell, Grimshaw \& Williams, 2010}
\end{frame}

% The sweet spot of control (schematic, from the book's argument)
\begin{frame}{Վերահսկողության «ոսկե միջինը»}
  \begin{columns}[c]
    \begin{column}{0.56\textwidth}
      \centering
      \begin{tikzpicture}
        \draw[-{Stealth[length=2.4mm]}, plum, line width=1pt] (0,0) -- (6.4,0);
        \draw[-{Stealth[length=2.4mm]}, plum, line width=1pt] (0,0) -- (0,3.4);
        \node[tag, font=\small, rotate=90, anchor=south] at (-0.05,1.7) {«զվարճալի վախ»};
        \node[tag, font=\small] at (3.2,-1.35) {վերահսկողություն $\rightarrow$};
        \draw[amOrange, line width=2.2pt] (0.2,0.2) .. controls (1.6,0.4) and (1.9,2.9) .. (3.0,2.9) .. controls (4.2,2.9) and (4.8,0.6) .. (6.1,0.4);
        \node[visible on=<2->] at (0.7,1.1) {\bigemoji[0.6cm]{skull}};
        \node[tag, text width=1.9cm, visible on=<2->] at (0.9,-0.6) {առանց պարաշյուտի};
        \node[visible on=<3->] at (3.0,3.35) {\bigemoji[0.6cm]{rollercoaster}};
        \node[tag, text width=2.0cm, visible on=<3->] at (3.2,-0.6) {պարաշյուտ, ամերիկյան բլուրներ};
        \node[visible on=<4->] at (5.6,1.1) {\bigemoji[0.6cm]{videogame}};
        \node[tag, text width=1.4cm, visible on=<4->] at (5.6,-0.6) {տեսախաղ};
      \end{tikzpicture}\par\smallskip
      {\scriptsize Սխեմատիկ՝ ըստ գրքի դատողության}
    \end{column}
    \begin{column}{0.4\textwidth}
      \begin{itemize}
        \item Որքան \textbf{քիչ} ես վերահսկում, այնքան ուժեղ է սարսուռը
        \item<2-> Բայց կա սահման. պարաշյուտով ցատկը զվարճալի է, առանց պարաշյուտի՝ ոչ
        \item<3-> Պետք է իրական ռիսկ, բայց նաև ելքի վրա \textbf{ազդելու} հնարավորություն
        \item<5-> Ավտովթարից փրկվածները ուրախ են, բայց կրկնել չեն ուզում
      \end{itemize}
    \end{column}
  \end{columns}
\end{frame}

% Counterfactual thinking
\begin{frame}{«Իսկ եթե...»}
  \begin{columns}[c]
    \begin{column}{0.58\textwidth}
      \begin{itemize}
        \item Ճանապարհն անցնելիս մեքենան քեզ համարյա դիպավ
        \item<2-> Ֆիզիկապես ոչինչ չի փոխվել. ողջ ու առողջ ես
        \item<3-> Բայց օրերով մտածում ես, թե ինչ \textbf{կլիներ}, եթե դիպչեր
        \item<4-> Սա \textbf{կոնտրֆակտուալ մտածողություն} է (counterfactual thinking). չեղած դեպքերի վատ ելքերի մասին մտածելը
        \item<5-> Ուղեղը սիրում է կենտրոնանալ սպառնալիքի վրա՝ անցյալում, ներկայում թե ապագայում
      \end{itemize}
    \end{column}
    \begin{column}{0.36\textwidth}
      \centering
      \begin{tikzpicture}
        \node at (0,-0.6) {\mascot[2.3cm]{thinking}};
        \node[think, text width=2.4cm, font=\mediumfont\small, callout absolute pointer={(0.3,0.45)}, visible on=<3->] at (0.6,1.9) {Իսկ եթե դիպչե՞ր...};
        \node at (-1.7,1.3) {\bigemoji[0.9cm]{car}};
      \end{tikzpicture}
    \end{column}
  \end{columns}
\end{frame}

% Sensation seeking: myth or fact
\begin{frame}{Առասպել, թե՞ փաստ \#2}
  \mythorfact{Սուր զգացողություններ փնտրողներն ավելի լավ են գնահատում ռիսկը}{ԱՌԱՍՊԵԼ}{pinkdark}{%
    \begin{card}[colback=brainpink!25!cardwhite, top=1mm, bottom=1mm]
      \textbf{Սուր զգացողությունների որոնում} (sensation seeking)՝ նորի ու ինտենսիվի ձգտում՝ անգամ ռիսկի գնով։ Վնասված դահուկորդներից շատերն այդպիսին էին։
    \end{card}}
  \source{Neary \& Zuckerman, 1976; Bouter et al., Pers. Individ. Differ., 1988}
\end{frame}

% Meme: the skier
\begin{frame}{Սայթաքուն լանջ}
  \begin{columns}[c]
    \begin{column}{0.5\textwidth}
      \begin{itemize}
        \item Սարսուռի ձգտումը մղում է քայլերի, որոնք դուրս են գալիս վերահսկողությունից
        \item Ռիսկ փնտրելու ցանկությունը խանգարում է \textbf{տեսնել} ռիսկը
      \end{itemize}
      \vspace{2mm}
      \onslide<2->{%
      \begin{bookquote}
        Գուցե ամեն ինչ աստիճանաբար է լինում. դասական «սայթաքուն լանջ»։ Դահուկորդների համար շատ տեղին տերմին է։
      \end{bookquote}}
    \end{column}
    \begin{column}{0.44\textwidth}
      \centering
      \begin{meme}[5.2cm]{surprised_pikachu.jpg}
        \memecap[font=\blackfont\scriptsize, fill=none, draw=none]{0.5}{0.86}{Ես. սահում եմ իմ հնարավորություններից\\ շատ ավելի արագ}
        \memecap[font=\blackfont\scriptsize, fill=none, draw=none]{0.5}{0.68}{Իմ ծունկը. կոտրվում է}
        \memecap{0.75}{0.12}{ես}
      \end{meme}
    \end{column}
  \end{columns}
\end{frame}

% DRD4
\begin{frame}{Գենե՞րը. DRD4}
  \begin{columns}[c]
    \begin{column}{0.46\textwidth}
      \begin{itemize}
        \item DRD4 գենը կոդավորում է դոֆամինի ընկալիչ։ Սուր զգացողություններ փնտրողների մոտ այն երբեմն տարբեր է
        \item<2-> Գուցե պարգևի ուղին \textbf{այլ կերպ} է աշխատում
      \end{itemize}
      \onslide<4->{%
      \begin{scienceupdate}
        Մետավերլուծությամբ կապը շատ թույլ է կամ չկա (առավելագույնը 3\%)։
      \end{scienceupdate}}
    \end{column}
    \begin{column}{0.5\textwidth}
      \centering
      \begin{tikzpicture}[ar/.style={-{Stealth[length=2.2mm]}, plum, line width=1pt}]
        \node[pbox=cardwhite, text width=2.2cm, visible on=<2->] (g) at (0,0) {\emoji{dna}\ DRD4};
        \node[pbox=mint!40!cardwhite, text width=2.9cm, visible on=<2->] (a) at (-1.75,-1.55) {ուղին \textbf{ավելի} ակտիվ է. ամեն ինչ ուժեղ է զգացվում};
        \node[pbox=lavender!40!cardwhite, text width=2.9cm, visible on=<3->] (b) at (1.75,-1.55) {ուղին \textbf{պակաս} ակտիվ է. ավելի ուժեղ խթան է պետք};
        \draw[ar, visible on=<2->] (g) -- (a); \draw[ar, visible on=<3->] (g) -- (b);
        \node[pbox=peach!45!cardwhite, text width=5cm, visible on=<3->] (c) at (0,-3.15) {Երկու դեպքում էլ՝ ավելի շատ սարսուռ փնտրել};
        \draw[ar, visible on=<3->] (a) -- (c); \draw[ar, visible on=<3->] (b) -- (c);
      \end{tikzpicture}
    \end{column}
  \end{columns}
  \only<1-3>{\source{Zuckerman, Genetics of sensation seeking (գրքից)}}%
  \only<4->{\source{Munafò et al., Biol. Psychiatry, 2008}}
\end{frame}

% The novelty detector in the hippocampus
\begin{frame}{«Սա արդեն տեսե՞լ ենք»}
  \begin{columns}[T]
    \begin{column}{0.54\textwidth}
      \begin{itemize}
        \item Սառա Մարտինն ու գործընկերները (2007) սկանավորեցին տասնյակ մարդկանց ուղեղ
        \item<2-> Որքան ուժեղ էր նորության ձգտումը, այնքան \textbf{մեծ} էր աջ հիպոկամպի առաջնային մասը
        \item<3-> Դա նորույթի դետեկտոր է. «Սա արդեն տեսե՞լ ենք»
      \end{itemize}
    \end{column}
    \begin{column}{0.42\textwidth}
      \onslide<4->{%
      {\blackfont\small Ինչու՞ է մեծ. երկու վարկած}\par\smallskip
      \begin{card}[colback=mint!35!cardwhite]
        \textbf{Ա}. Այնքան նոր բան է տեսել, որ «մեծացել է»
      \end{card}
      \begin{card}[colback=lavender!35!cardwhite]
        \textbf{Բ}. Չափից զարգացած է. «նոր» համարվելու համար շատ անսովոր բան է պետք
      \end{card}}
    \end{column}
  \end{columns}
  \onslide<5->{%
  \begin{card}[colback=peach!40!cardwhite]
    \centering Ամեն ուղեղ տարբեր է. վախենալու բան չկա, նույնիսկ եթե սիրում ես վախենալ
  \end{card}}
  \source{Martin et al., Neuropsychologia, 2007}
\end{frame}
